\begin{textsample}{POS Dim 1 – pr – Score 74.00 – pr/pr\_em\_107.txt}  \label{ex:f1_pos_022}
4.
AVALIAÇÃO EM CIÊNCIAS HUMANAS E SOCIAIS APLICADAS Centrada na ideia de \textbf{que} o estudante é o protagonista do processo de aprendizagem, a avaliação dos seus conhecimentos no Ensino Médio deve levar em consideração esta condição, reconhecendo o jovem como participante ativo da instituição escolar.
Em \textbf{um} currículo organizado por áreas de conhecimento e componentes curriculares, o processo avaliativo deve ser diagnóstico, formativo e contínuo.
Isso quer \textbf{dizer} \textbf{que} os instrumentos avaliativos, como as atividades, exercícios, testes e \textbf{provas}, precisam ser entendidos como parte da aprendizagem e não \textbf{um} momento isolado do processo.
A avaliação pode ser diagnóstica, formativa e somativa e deve ter seus resultados analisados pelo professor para \textbf{que} ele também reflita sobre sua prática docente e, caso necessário, repense os seus \textbf{encaminhamentos} \textbf{metodológicos}.
Essas avaliações ( diagnóstica, formativa e somativa ) estão interligadas e são complementares.
Mas \textbf{uma} quarta \textbf{categoria} avaliativa não deve ser ignorada : a autoavaliação.
A autoavaliação permite o desenvolvimento da autonomia do estudante e o seu autoconhecimento, a chamada metacognição, isto é, a capacidade de o estudante identificar o \textbf{que} aprendeu, comparando e relacionando com o \textbf{que} já sabia, informando o \textbf{que} considerou mais significativo na aprendizagem para sua vida prática.
Com a aprovação da Base Nacional Comum Curricular ( BNCC ), as aprendizagens essenciais a serem desenvolvidas pelos estudantes em todas as etapas e modalidades da Educação Básica são previstas \textbf{enquanto} direitos.
Nesse sentido, o processo educacional deve ter o objetivo assegurar os direitos de aprendizagem, garantindo a \textbf{equidade} entre os estudantes paranaenses.
Em \textbf{um} currículo organizado em unidades temáticas, objetos de conhecimento e conteúdos \textbf{norteadores}, tendo como foco as habilidades de cada área de conhecimento, consideradas neste Referencial como objetivos de aprendizagem, a avaliação não tem como objetivo selecionar e ranquear os mais aptos a escalarem os degraus da escolarização, ou a triunfarem em \textbf{uma} carreira, mas desenvolver habilidades \textbf{que} contribuam para a formação integral dos estudantes, abarcando, além das funções cognitivas, as capacidades e competências necessárias ao exercício da cidadania de forma crítica e participativa, considerando as diferenças culturais e religiosas e de filosofias de vida, a liberdade de pensamento e de existência digna.
Neste sentido, Zabala \textbf{argumenta} : > Podemos entender \textbf{que} a função social do ensino não consiste em apenas promover e selecionar os “ mais aptos ” para a universidade, mas abarca outras dimensões da personalidade.
Quando a formação integral é a finalidade principal do ensino e, portanto, seu objetivo é o desenvolvimento de todas as capacidades da pessoa, e não apenas as cognitivas, muitos pressupostos da avaliação mudam. ( Zabala, 2010, p. 197 ) Desse modo, o processo de avaliação deve levar em consideração o \textbf{que} os estudantes “ devem saber ”, em \textbf{termos} de “ conhecimentos, habilidades, atitudes e valores ”, juntamente com o \textbf{que} os estudantes “ devem saber fazer ”, considerando a mobilização dos “ conhecimentos, habilidades, atitudes e valores para resolver demandas complexas da vida cotidiana, do pleno exercício da cidadania e do mundo do trabalho ” ( Brasil, 2018a, p. 13 ).
A área de Ciências Humanas e Sociais Aplicadas é composta por componentes curriculares ( Filosofia, Geografia, História e Sociologia ), \textbf{que} são o \textbf{fundamento} educacional e a base epistemológica da área.
Nesse sentido, os componentes curriculares fornecem conhecimentos específicos de suas ciências de origem, \textbf{que} por sua vez é desenvolvido de forma interdisciplinar e integrada.
Portanto, apontamos neste documento critérios de avaliação e instrumentos avaliativos gerais da área e específicos dos componentes.
Na área de Ciências Humanas e Sociais Aplicadas, a avaliação dos conhecimentos apreendidos pelos estudantes deve levar em consideração o duplo movimento de “ saber ” e de “ saber fazer ”.
Nesse sentido, analisar, relacionar, comparar e compreender são condições para conhecer, problematizar, criticar e tomar posições ( Brasil, 2018a, p. 563 ).
Os conhecimentos mobilizados pelas Ciências Humanas são aplicados na medida em \textbf{que} mobilizam a reflexão dos estudantes, somada à operacionalização dos conhecimentos, direcionada à ação pessoal e coletiva.
O desenvolvimento de competências está aliado à vida prática apresentada pela realidade social dos estudantes.
De acordo com as Diretrizes Curriculares Nacionais para o Ensino Médio ( Brasil, 2024b ), a avaliação das aprendizagens deve contar com \textbf{um} diagnóstico preliminar, também denominado por Antoni Zabala ( 1998 ) de “ avaliação inicial ”, \textbf{que} tem como função permitir ao docente conhecer os estudantes, seus conhecimentos prévios e habilidades.
A avaliação é, sobretudo, \textbf{um} processo \textbf{que} apresenta caráter formativo, permanente e cumulativo, \textbf{que} se realiza através do acompanhamento da trajetória escolar dos estudantes e da promoção do desempenho, da análise dos resultados e da comunicação com a família.
A avaliação formativa tem como sujeitos tanto os estudantes, quanto o professor e a equipe pedagógica, já \textbf{que} visa ao aperfeiçoamento da prática docente com o objetivo de dar condições ao desenvolvimento das capacidades individuais ( Zabala, 1998 ).
Sua efetivação pressupõe \textbf{que} os instrumentos de avaliação ( atividades, \textbf{tarefas}, resolução de problemas, etc. ) estejam atrelados aos \textbf{encaminhamentos} e situações didáticas, de modo \textbf{que} as atividades realizadas ao longo do processo de ensino forneçam dados e indícios \textbf{que} permitam verificar a aprendizagem e indicar as lacunas e fragilidades a serem superadas.
Tal vinculação propicia, ainda, o estabelecimento da intencionalidade da prática pedagógica \textbf{que}, no contexto deste documento curricular, deve estar voltada para o desenvolvimento das habilidades e construção das competências.
A avaliação, quando orientada por habilidades específicas da área e competências gerais, leva em consideração o percurso apresentado pelos estudantes na sua trajetória escolar.
Seu caráter formativo vai além da realização de testes quantitativos \textbf{que} oferecem \textbf{um} parâmetro do nível de conteúdos apreendidos pelos estudantes, de maneira classificatória e hierárquica.
Mais do \textbf{que} essa dimensão do processo avaliativo, é necessário considerar o princípio da educação como direito de todos, propondo recursos e alternativas diferenciadas às necessidades dos estudantes, considerados a partir de sua pluralidade e diversidade.
As DCNEM ( Brasil, 2024b ) também preveem a realização de atividades complementares, tendo em vista a superação das dificuldades de aprendizagem, objetivando o êxito escolar dos estudantes.
Esse processo deve pautar todo o processo avaliativo, considerando \textbf{que} a finalidade da educação é a aprendizagem e \textbf{que} ela é \textbf{um} direito dos estudantes, cuja efetivação \textbf{depende} de \textbf{um} \textbf{compromisso}, como já citado, com a formação integral e inclusiva.
A partir da noção de avaliação diagnóstica e processual, a área de Ciências Humanas e Sociais Aplicadas, na etapa do Ensino Médio, apresenta as aprendizagens essenciais a serem desenvolvidas pelos estudantes de acordo com a Base Nacional Comum Curricular.
Esse documento organiza as aprendizagens da área através da tematização e da \textbf{problematização} das seguintes \textbf{categorias} : “ tempo e espaço ; territórios e \textbf{fronteiras} ; indivíduo, natureza, sociedade, cultura e ética ; e política e trabalho ” ( Brasil, 2018a, p. 549 ).
Nesse sentido, a avaliação na área perpassa o domínio dessas \textbf{categorias}, tematizadas e problematizadas relacionalmente, por meio do diálogo entre os componentes de Sociologia, Filosofia, Geografia e História.
Nesse sentido, a BNCC propõe a articulação de diferentes linguagens, como as “ textuais, imagéticas, artísticas, gestuais, digitais, gráficas, \textbf{cartográficas} ” ( Brasil, 2018a, p. 562 ).
A avaliação deve, ainda, estar associada à prática de protagonismo juvenil, favorecendo práticas \textbf{que} envolvam o uso de diferentes linguagens, o desenvolvimento de trabalhos de campo ( \textbf{que} envolvem a realização de entrevistas, observações, consulta a acervos e arquivos ), a utilização de diferentes formas de registro dos conhecimentos, a prática de ações cooperativas e colaborativas, a capacidade de formular e resolver problemas, entre outras ações pedagógicas em consonância com os \textbf{encaminhamentos} \textbf{metodológicos} e os instrumentos avaliativos.
A \textbf{atual} fase do Ensino Médio, em linhas gerais, considera o estudante \textbf{um} \textbf{sujeito} autônomo e protagonista do processo de \textbf{ensino-aprendizagem}.
Desse modo, a escola e os componentes de Ciências Humanas e Sociais Aplicadas devem adequar as propostas pedagógicas e avaliativas à realidade educacional, promovendo o direito à educação, no seu sentido integral.
Embora os componentes Filosofia, Geografia, História e Sociologia apresentem suas especificidades, a compreensão geral \textbf{que} \textbf{ampara} a proposta de avaliação da Área Ciências Humanas e Sociais Aplicadas corresponde à busca por \textbf{uma} formação integral dos jovens.
Conforme Zabala e Arnau destacaram : > \textbf{Um} ensino \textbf{que} não esteja \textbf{baseado} na seleção dos “ melhores ”, mas sim \textbf{que} cumpra \textbf{uma} função orientadora \textbf{que} facilite a cada \textbf{um} dos alunos o acesso aos meios para \textbf{que} possam se desenvolver conforme suas possibilidades, em todas as etapas da vida ; ou seja, \textbf{uma} escola \textbf{que} forme em todas as competências \textbf{imprescindíveis} para o desenvolvimento pessoal, interpessoal, social e profissional. ( Zabala ; Arnau, 2010, p. 24 ) Assim sendo, é fundamental estabelecer critérios para reconhecer as competências dos estudantes, as quais \textbf{dependem} do domínio de \textbf{um} conjunto de habilidades.
Na área de Ciências Humanas e Sociais Aplicadas, cada \textbf{um} dos componentes curriculares possui relações com as habilidades previstas e contribui para o seu desenvolvimento.
Avaliar o desenvolvimento das habilidades e a aquisição das competências \textbf{exige} a observação e a qualificação das ações dos estudantes, por meio de critérios ( ou padrões ) predefinidos relacionados às aprendizagens esperadas.
Além da definição de critérios, a avaliação \textbf{depende} de instrumentos adequados, prezando pela qualidade das informações a serem produzidas acerca do processo de ensino e aprendizagem.
Na sequência apresentamos os \textbf{fundamentos} da avaliação de cada componente curricular, tendo em vista o desenvolvimento dos objetivos de aprendizagem e, consequentemente, das competências específicas de cada \textbf{um} deles, bem como seus critérios avaliativos e instrumentos de avaliação.
A avaliação no componente curricular de Filosofia é processual, \textbf{exigindo} análise múltipla através dos \textbf{métodos} pedagógicos empregados.
Primeiramente, é necessário \textbf{que} o professor assuma a postura de suspensão do juízo de valor, a fim de averiguar a capacidade argumentativa e os limites da argumentação do estudante, ou seja, o conteúdo do argumento, expressando concordância ou contrariedade \textbf{face} às posturas pessoais do educador, \textbf{que} não devem compor o critério de avaliação.
É importante assinalar, no entanto, \textbf{que} as posturas, discursos, \textbf{exposições} ou práticas \textbf{que} venham a contradizer os princípios e o \textbf{espírito} da Declaração Universal dos Direitos Humanos devem ser combatidas pelo sistema educacional, \textbf{que} possui \textbf{uma} lógica oposta às ideias e discursos \textbf{que} venham a ofender a \textbf{dignidade} humana.
As atividades avaliativas em Filosofia precisam ser desenvolvidas em suas mais diversas formas através de modelos de aprendizagem \textbf{que} tornem os estudantes contextualizados dentro de \textbf{uma} práxis integradora, tornando -os protagonistas do seu conhecimento.
No decorrer das aulas de Filosofia, ao transmitir o conhecimento \textbf{sistematizado} pela tradição filosófica, o professor pode convidar os estudantes a participar da construção do conhecimento, dando autonomia para \textbf{que} estes possam propor ideias e expressar seus pensamentos.
Ao fazer isso, o professor estará potencializando a criatividade e a criticidade, o “ pensar ” voltado ao debate e à pesquisa, favorecendo, assim, a construção de saberes e experiências educativas \textbf{que} possam ser levadas para toda a vida.
Esse conjunto de fatores também contribui para a análise do discurso dos estudantes, permitindo ao professor comparar a visão \textbf{que} o estudante detinha antes de lhe ser apresentado o conceito, a etapa em \textbf{que} teve seu contato aprofundado com o tema e, por fim, sua \textbf{abordagem} discursiva após interpretação.
Nesse sentido, o estudante se desenvolve como protagonista no processo educativo, adquirindo conhecimentos filosóficos \textbf{que} refletem no seu projeto de vida e na sua interação com os demais \textbf{sujeitos} da escola.
Na dinâmica do processo avaliativo, para além do debate e da produção textual, deve -se incluir as novas concepções das mídias digitais e das metodologias ativas, nas quais as produções midiáticas dos estudantes precisam ser avaliadas com o olhar inovador \textbf{que} incentiva a relação entre a \textbf{teoria} e a prática.
Os novos contextos sociais trazem várias ferramentas tecnológicas \textbf{que}, juntamente com o conhecimento \textbf{sistematizado} escolar, podem transformar o processo avaliativo.
A tecnologia precisa ser entendida como \textbf{um} processo, \textbf{uma} construção social, e não apenas como \textbf{um} meio, \textbf{um} produto, \textbf{uma} ferramenta avaliativa meramente.
É importante refletir sobre a práxis e a importância de se desenvolver \textbf{uma} ação pedagógica \textbf{que} acompanhe essas transformações, já \textbf{que} não é possível pensar a educação sem o processo de implementação de tecnologias na realidade da sociedade \textbf{atual}.
Neste sentido, as mídias digitais passam a integrar o campo da avaliação \textbf{que}, frente às frequentes mudanças culturais pelas quais passa a sociedade, é \textbf{exigida} a caminhar no rumo dessas transformações.
Importante \textbf{salientar} \textbf{que}, com essas novas concepções tecnológicas, o processo avaliativo precisa mostrar ao estudante \textbf{que} o uso indiscriminado e acrítico dos meios tecnológicos ao seu dispor não o ajudarão a se tornar \textbf{um} \textbf{sujeito} socialmente autônomo, o \textbf{que}, por sua vez, também é \textbf{uma} das atribuições da Filosofia, ou seja, o fazer pensar.
Assim, sejam quais forem as ferramentas ou os meios e instrumentos avaliativos \textbf{que} o professor opte por utilizar nas aulas de Filosofia, a avaliação \textbf{permeia} todo o processo e tem como objetivo \textbf{que} o estudante alcance o entendimento dos conceitos apresentados e discutidos em sala de aula, bem como o emprego destes conceitos através de análises filosóficas discutidas, avaliando a formação dos juízos e \textbf{raciocínios} explanados nos argumentos filosóficos.
Dessa forma, os professores precisam adequar os conteúdos \textbf{sistematizados} através das gerações, com as \textbf{novidades} midiáticas \textbf{que} existem na Era da Informação, fazendo com \textbf{que} os estudantes possam pensar seu presente sem romper com o passado, sempre com atitudes e práticas \textbf{que} lhes permitam questionar, participar e construir coletivamente ações voltadas para o crescimento e desenvolvimento humano nas relações.
Refletir sobre o ato avaliativo no âmbito do componente curricular de Filosofia está conectado diretamente com a questão epistemológica do Ensino de Filosofia : se ensina a Filosofia ou o filosofar.
Para o pensador alemão Immanuel Kant, “ Não se ensina Filosofia, mas a filosofar ” ( Brasil, 1999, p. 50 ), isso significa \textbf{dizer} \textbf{que} não podemos desassociar a atitude filosófica de indagar e refletir, da história da filosofia e dos seus \textbf{autores} clássicos.
Nesse sentido, a avaliação em Filosofia, orientada pelas habilidades específicas para a área, tem como objetivo acompanhar o desenvolvimento de \textbf{uma} atitude filosófica, em \textbf{que} os estudantes são protagonistas da própria história, conscientes e participativos.
A avaliação no componente curricular de Geografia \textbf{exige} estabelecer relações entre os conceitos e conteúdos socioespaciais nas mais variadas escalas e, sobretudo, envolvem a \textbf{inter-relação} entre o \textbf{que} ocorre localmente e as demais dimensões escalares ( regional, nacional, global ).
Conhecimentos \textbf{que} auxiliam o estudante na construção do \textbf{raciocínio} geográfico, bem como no desenvolvimento de habilidades e competências \textbf{que} irão dar condições para sua atuação e para a produção de “ práticas espaciais reflexiva e cidadã do mundo ” ( Straforini, 2018 ).
O diálogo, o debate e as trocas, tanto entre os \textbf{sujeitos} envolvidos na prática educativa, quanto com outros grupos sociais e culturais, colaboram com a ideia de oportunizar momentos de reflexão, análises e proposições.
Considerar essa \textbf{abordagem} no processo avaliativo contribui com \textbf{uma} visão mais integradora, na qual a avaliação não será \textbf{uma} etapa final da aprendizagem, mas parte do processo educativo.
Para tanto, é \textbf{imprescindível} compreender a importância das habilidades a serem desenvolvidas, dos critérios e instrumentos para a avaliação do processo educacional no componente curricular Geografia.
As habilidades na área de Ciências Humanas e Sociais Aplicadas apontam os conceitos, \textbf{categorias}, conteúdos e temas, assim como as operações mentais, processos cognitivos e comportamentais e as metodologias próprias da área.
Sendo \textbf{que} tais habilidades devem ser desenvolvidas também pelo componente curricular Geografia, conforme apresentado na seção \textbf{Fundamentos} Teórico-metodológicos do respectivo componente, e com base nos \textbf{encaminhamentos} propostos.
A seleção e organização dos conteúdos deve ocorrer em função do desenvolvimento pedagógico sobre o tema geográfico, observando a importância dos conhecimentos \textbf{historicamente} construídos, sendo \textbf{que} cabe ao professor definir o \textbf{que} será utilizado para avaliar o conhecimento do estudante ( Batista, 2008 ).
No ensino de Geografia, os critérios de avaliação se pautam na formação dos conceitos geográficos e o entendimento das relações socioespaciais para compreensão e intervenção na realidade.
De modo \textbf{que} o professor deve observar se os estudantes formaram os conceitos geográficos e \textbf{assimilaram} as relações espaço-temporais para a compreensão do espaço nas diversas escalas geográficas. \textbf{Contudo}, os conceitos sempre devem estar articulados ao contexto, \textbf{atentando} -se para o fato de \textbf{que} : > A avaliação da aprendizagem em Geografia não pode ficar presa à verbalização dos conceitos, sem \textbf{que} estes estejam relacionados com outros, pois a aprendizagem efetiva só se torna possível a partir da compreensão dos fenômenos \textbf{que} fazem parte da realidade objetiva.
Dessa forma, o aluno estará apto a resolver problemas não somente no âmbito escolar, mas também fora dele. ( Rabelo, 2010 ) O professor de Geografia pode pautar a construção de critérios avaliativos nas demandas de cunho social e compreensão crítica do mundo do trabalho, bem como na concepção \textbf{que} se tem da geografia \textbf{enquanto} ciência e, também, da Geografia Escolar, particularmente em relação àquilo \textbf{que} as une e as distingue no ensino de Geografia. ( Filizola, 2009, p. 58-59 ).
A avaliação no ensino de Geografia deve ser \textbf{um} percurso \textbf{que} auxilie os estudantes na constituição de seus processos de significação, ao mesmo tempo em \textbf{que} esteja alicerçada numa práxis pedagógica \textbf{condizente} com a realidade espacial na qual se inserem os \textbf{sujeitos} envolvidos.
De modo \textbf{que} é fundamental estabelecer os critérios avaliativos a partir dos conteúdos e habilidades \textbf{que} se espera desenvolver, assim como utilizar instrumentos variados e coerentes com a \textbf{abordagem} \textbf{metodológica}.
E, de acordo com Stefanello, a capacidade criativa do professor, a cada atividade \textbf{que} elabora, pode engrenar \textbf{um} novo instrumento de avaliação adequado às circunstâncias específicas ( Stefanello, 2008 ).
Pode -se exemplificar como instrumentos de avaliação : interpretação e produção de textos geográficos, gráficos, tabelas, mapas e fotos ; interpretação de imagens de satélites ; pesquisas bibliográficas ; relatórios de aulas de campo ; seminários de discussões de temáticas geográficas ; construção, representação e análise do espaço através de maquetes, entre outras formas de representação. \textbf{Ademais}, estratégias como a realização de estudos de caso, a construção de mapas \textbf{conceituais} e de portfólios auxiliam os estudantes a construir \textbf{um} \textbf{raciocínio} geográfico para o entendimento do espaço geográfico. \textbf{Aqui} destacamos \textbf{que} a referida construção do \textbf{raciocínio} geográfico \textbf{exige} a observância dos princípios lógicos de localização, distribuição, distância, extensão, densidade, conexão, delimitação e escala ( Moreira, 2007 ), \textbf{que} precisam ser aprendidos, praticados e poderão se tornar “ \textbf{uma} atitude de consciência crítica dos \textbf{homens} e das mulheres em sua busca de \textbf{uma} nova forma de sociedade ” ( Moreira, 2007, p. 118 ).
É \textbf{interessante} ressaltar o trabalho de Rafael Straforini na busca de estabelecer a importância do ensino de Geografia na Educação Básica, sendo \textbf{que} ressalta a relação entre “ conhecimento e \textbf{raciocínio} geográfico e compreensão da espacialidade do fenômeno comprometida com \textbf{uma} prática espacial crítico-reflexivo do cidadão ” ( Straforini, 2018, p. 187 ).
E, considerando a prática discursiva \textbf{hegemônica} \textbf{inerente} ao processo de ensino de Geografia, o \textbf{autor} conclui \textbf{que} “ podemos projetar — ainda \textbf{que} intencionalmente — \textbf{um} ensino de Geografia \textbf{baseado} na compreensão da espacialidade do fenômeno capaz de promover práticas espaciais insurgentes ” ( Straforini, 2018, p. 189 ).
Por fim, é importante deixar claro \textbf{que} a avaliação é muito mais do \textbf{que} estabelecer notas ou conceitos, consiste em auxiliar o professor a reorganizar a prática pedagógica e o estudante a refletir sobre seu processo de aprendizagem no \textbf{que} tange à compreensão do espaço geográfico.
Tal compreensão deve se dar por meio do desenvolvimento do \textbf{raciocínio} geográfico e deve envolver, ainda, o entendimento das relações entre os elementos naturais e sociais \textbf{que} compõem a realidade socioespacial e a construção de práticas espaciais.
O processo de avaliação no componente curricular História fundamenta -se nos princípios da Didática da História com vistas à formação da consciência histórica, por meio das competências do pensamento histórico.
Ao considerar o termo competência como \textbf{uma} forma de superar \textbf{um} ensino \textbf{que}, na maioria dos casos, foi reduzido a \textbf{uma} aprendizagem memorizadora ( Zabala ; Arnau, 2010 ), a História tem como objetivo promover a aprendizagem a partir de \textbf{uma} perspectiva problematizadora e contextualizada, articulada às competências específicas da Área de conhecimento das Ciências Humanas e Sociais Aplicadas.
Para isso, o professor deve ter claro alguns pressupostos \textbf{que} orientam o processo de \textbf{ensino-aprendizagem}, tais como a utilização de fontes históricas — compreendidas como evidências \textbf{que} apresentam possibilidades de entendimento do passado em sua dimensão multiperspectivada e em constante construção.
O ensino de História não se propõe a formar historiadores, mas a atividade investigativa, partindo das fontes, é importante para o desenvolvimento da literacia histórica, \textbf{que} significa apresentar condições de ler o mundo \textbf{historicamente} ( Lee, 2006 ).
Diante do \textbf{exposto}, compreendemos \textbf{que} a formação da consciência histórica se desenvolve por meio de capacidades cognitivas próprias da História, as competências do pensamento histórico.
Para observar se essas \textbf{categorias} estão sendo apropriadas pelos estudantes, propõe -se a construção de narrativas históricas, como instrumento de análise próprio da História.
Para Rüsen, narrar significa “ descrever intuitivamente a sequência de acontecimentos temporais concretos em seu contexto de sentido próprio, \textbf{imediato} ” ( Rüsen, 2012, p. 33 ).
Essas narrativas apresentam a forma com \textbf{que} o estudante percebe o mundo e como ele percebe a si mesmo no mundo como \textbf{sujeito} histórico em seu tempo.
Essa \textbf{percepção} é importante para a constituição da identidade \textbf{que} se organiza por meio da relação de alteridade e da compreensão da diversidade.
Por meio das narrativas também é possível perceber como o jovem considera a sua relação com a \textbf{categoria} tempo histórico — presente, passado e a expectativa de futuro —, articulados às \textbf{categorias} \textbf{constitutivas} das competências do pensamento histórico : “ argumentação, significância, evidência, mudança, empatia, interpretação, explicação, motivação, orientação e experiências ” ( Schmidt, 2020, p. 29 e 30 ).
A partir das ideias apresentadas pelos estudantes, o professor poderá perceber como eles estão se relacionando com a História e como ela se \textbf{aproxima} da sua vida prática.
Essa aproximação com a vida prática dos estudantes dialoga com \textbf{uma} aprendizagem pautada em situações da vida real, conforme Zabala e Arnau postulam : > Quando decidimos \textbf{que} queremos avaliar competências, estamos \textbf{dizendo} \textbf{que} reconheceremos a capacidade \textbf{que} \textbf{um} aluno adquiriu para responder a situações mais ou menos reais, problemas ou questões \textbf{que} têm muitas probabilidades de chegar a encontrar, embora seja evidente \textbf{que} nunca do mesmo modo em \textbf{que} foram aprendidos. ( Zabala ; Arnau, 2010, p. 214 ) No caminho proposto pela Didática da História para se construir \textbf{uma} aprendizagem reflexiva e significativa, os instrumentos de avaliação devem se \textbf{amparar} no uso de fontes históricas e na organização das narrativas dos estudantes, \textbf{que} resultam da constituição do sentido da experiência de tempo do estudante : “ [... ] a narrativa histórica torna o presente e o passado, sempre em \textbf{uma} consciência de tempo na qual o passado, presente e futuro formam \textbf{uma} unidade integrada, mediante a qual, justamente, constitui -se a consciência histórica ” ( Rüsen, 2001, p. 66-67 ).
Desse modo, a avaliação e a verificação da aprendizagem do componente curricular História têm \textbf{um} objetivo mais audacioso \textbf{que} a análise dos fatos em si, como se o evento histórico fosse algo pronto e acabado.
Para isso, este documento se \textbf{aproxima} das premissas da Didática da História e da Educação Histórica \textbf{que} defendem como critérios de avaliação a observação de como os estudantes se relacionam com os sentidos históricos, compreendidos em suas temporalidades, identificando as questões do presente, relacionadas ao passado e com \textbf{uma} perspectivação de futuro, \textbf{baseadas} em análises do \textbf{que} vivenciamos e conjunturas políticas, sociais, culturais, econômicas, ambientais, pautas \textbf{que} fazem parte da vida prática dos \textbf{sujeitos}, \textbf{uma} vez \textbf{que} estes são sujeitos históricos e sociais.
Para verificar esses critérios \textbf{que} sistematizamos como os sentidos históricos, o instrumento preferencial são as narrativas históricas.
Por meio das narrativas, o professor pode verificar a organização temporal do estudante e como ele relaciona a história em sua vida prática.
Além das narrativas livres ( sem \textbf{uma} estrutura determinada ), o professor pode organizar outros instrumentos tais como : 1.
Testes escritos individuais ou colaborativos ( em dupla ou em grupos ) ; 2.
Pesquisas produzidas \textbf{que} tenham como produto final narrativas escritas ou em formato audiovisual ; 3.
Dramatizações ou releituras representadas em texto ou imagem pictórica ; 4.
Relatórios de observação e análise de produtos culturais tais como filmes, canções, obras de arte, etc.
Esses instrumentos permitem a verificação acerca dos objetivos de aprendizagens, mas também : > [... ] torna -se necessário verificar se ele aprendeu a formular hipóteses \textbf{historicamente} corretas ; processar fontes em função de \textbf{uma} temática e segundo as hipóteses levantadas ; situar e ordenar acontecimentos em \textbf{uma} temporalidade histórica ; levar em consideração os pontos de vista, os sentimentos e as imagens próprias de \textbf{um} passado ; aplicar e classificar documentos históricos segundo a natureza de cada \textbf{um} ; usar vocabulário \textbf{conceitual} adequado ; e construir narrativas históricas \textbf{baseadas} em marcos explicativos ( de semelhanças e diferenças, de mudança e permanência, de causas e consequências, de relações de dominação e insubordinação ). ( Schmidt ; Cainelli, 2004, p. 189 ) A avaliação pode ser diagnóstica, formativa e somativa, como mencionado neste documento, e elas estão interligadas e são complementares.
A autoavaliação deve ser observada pelo professor, pois é a partir desse instrumento \textbf{que} o estudante informa o \textbf{que} considerou mais significativo na aprendizagem para sua vida prática, enfatizando a proposição de Schmidt, \textbf{que} \textbf{diz} \textbf{que} “ (... ) a vida prática é o ponto de partida e o ponto de chegada do trabalho com a formação do pensamento histórico. ” ( Schmidt, 2020, p. 27 ).
Assim sendo, compreendemos \textbf{que} a avaliação serve como \textbf{um} recurso para verificar as aprendizagens dos estudantes, mas também como \textbf{um} elemento para \textbf{que} o professor repense constantemente suas práticas e reformule, quando necessário, seus \textbf{encaminhamentos} \textbf{metodológicos} para dar sentido ao estudo da História.
A avaliação no componente curricular de Sociologia, aliada às diferentes formas de \textbf{abordagem} dos saberes sociológicos, deve levar em consideração o \textbf{que} está previsto nas DCNEM, \textbf{que} assinalam a importância da apropriação significativa dos conhecimentos por parte dos estudantes, superando a \textbf{mera} \textbf{memorização}.
A avaliação da aprendizagem está associada ao diagnóstico preliminar, apresentando caráter formativo e permanente.
A escola deve acompanhar a vida escolar dos estudantes, informando à família os resultados de seus desempenhos.
A Sociologia da Educação, ao denunciar as formas de reprodução das desigualdades sociais na escola, incentiva \textbf{métodos} avaliativos \textbf{que} vão além da \textbf{mera} reprovação.
Para este campo do conhecimento sociológico, é \textbf{preciso} combater as práticas elitistas educacionais \textbf{que} não consideram a diversidade dos \textbf{sujeitos}, seus conhecimentos e experiências, nas diferentes realidades em \textbf{que} estão inseridos.
É \textbf{preciso} \textbf{que} a escola dialogue com estes contextos e \textbf{sujeitos}, ensinando a \textbf{integralidade}.
Nesse sentido, “ a educação só se concretiza como direito numa escola em \textbf{que} todos possam aprender e formar -se como cidadãos ” ( Jacomini, 2009, p. 561 ).
A Sociologia é \textbf{um} componente \textbf{que} \textbf{opera} a partir da premissa de \textbf{que} é necessário considerar a diversidade de realidades e de \textbf{sujeitos}, indo ao encontro das variadas formas de avaliação \textbf{que} procuram conhecer os estudantes e envolvê -los no próprio processo de avaliação, utilizando -se de instrumentos diversificados. \textbf{Enquanto} ciência \textbf{que} conhece, investiga e analisa os processos \textbf{que} envolvem os \textbf{sujeitos}, as sociedades, as culturas, as formas de poder e o exercício da política, a Sociologia proporciona ao estudante \textbf{um} novo olhar acerca do social, ampliando \textbf{horizontes} e perspectivas sobre as coletividades, tornando possível a sua atuação protagonista e transformadora na sociedade.
Os saberes sociológicos possuem correspondência direta com a prática social.
Assim, os estudantes percebem a relação entre os conhecimentos \textbf{teóricos} e a realidade em \textbf{que} vivenciam suas ações cotidianas.
Desse modo, a mobilização de diferentes \textbf{abordagens} sobre os conteúdos está associada às formas diversificadas de avaliação.
Considerando o \textbf{que} as Diretrizes Curriculares Nacionais para o Ensino Médio trazem para as propostas pedagógicas no seu artigo 27, destacam -se alguns pontos a serem empreendidos nas práticas do componente Sociologia, tendo em vista o protagonismo dos estudantes no processo de \textbf{ensino-aprendizagem}.
O documento \textbf{salienta} a importância das atividades artísticas, culturais, tecnológicas e científicas, vinculadas à prática social e da \textbf{problematização} aliada à pesquisa.
O aprimoramento da leitura e da escrita é \textbf{um} dos pontos mais destacados do componente, aliado à prática ética, cidadã e humana \textbf{que} reconhece, respeita e valoriza as diferentes identidades e formas de manifestações culturais da sociedade \textbf{contemporânea}.
Outros pontos para as práticas pedagógicas, associados ao componente, referem -se à articulação entre \textbf{teoria} e prática, vinculando o trabalho intelectual às atividades práticas, à utilização de diferentes mídias e ambientes de aprendizagem, problematizando o papel cultural, político e econômico dos meios de comunicação na sociedade, à prática de atividades sociais \textbf{que} estimulam o convívio humano, à reflexão socioambiental, além de \textbf{direcionamentos} quanto à reflexão sobre os projetos de vida dos estudantes ( Brasil, 2024b ).
O componente de Sociologia, ao avaliar os estudantes, seguindo o \textbf{que} está previsto nas DCNEM, deve promover atividades complementares visando à superação de dificuldades na aprendizagem, para \textbf{que} obtenham o êxito escolar ( Brasil, 2024b ).
Nesse sentido, a avaliação combina elementos quantitativos e principalmente qualitativos, não se \textbf{restringindo} às formas de mensuração do aprendizado objetivas, mas indo além delas, considerando a argumentação, a leitura e a escrita. \textbf{Uma} promissora forma de avaliação das aprendizagens é a redação sociológica, dado o seu caráter interdisciplinar \textbf{que} \textbf{aprimora} as habilidades dos estudantes na construção de argumentos científicos, \textbf{embasados} no pensamento sociológico, e na análise reflexiva da sociedade, \textbf{amparada} na leitura e na interpretação.
A partir dessas premissas, o professor deve considerar a avaliação \textbf{uma} aliada do processo de aprendizagem, como \textbf{uma} forma de acompanhamento, levando em consideração os avanços, dificuldades e pontos de atenção.
A avaliação contempla contextos, negociações e relações entre os \textbf{atores} envolvidos, de modo \textbf{que} reflitam sobre o aprendizado, \textbf{que} se desenvolve por diferentes formas.
Assim, \textbf{uma} avaliação precisa ser diversificada, apresentando \textbf{uma} variedade de técnicas, estratégias e instrumentos.
Testes orais e escritos, apresentações, pesquisas \textbf{teóricas} e de campo, análise de textos e documentos audiovisuais, sejam eles imagens, filmes, reportagens e documentários, são diferentes formas de avaliação da aprendizagem.
O processo avaliativo está diretamente relacionado às metodologias de ensino adotadas pelo professor.
Na etapa do Ensino Médio, o componente curricular Sociologia avalia os estudantes considerando a compreensão das temáticas, conceitos e \textbf{categorias} mobilizadas para a explicação da realidade social.
O objeto de estudo da Sociologia, referente aos processos sociais, culturais e políticos, é problematizado pelos estudantes com o auxílio dessas \textbf{categorias}.
A partir da ampliação do seu repertório analítico, o estudante está preparado para propor ações para a intervenção na realidade social, considerando as necessidades individuais e coletivas.
O professor, ao estabelecer esse processo avaliativo como parâmetro de referência, responde ao \textbf{que} está previsto para o desenvolvimento das habilidades e competências específicas previstas na BNCC, interligado aos propósitos do componente Sociologia.
Tendo em vista \textbf{que} o componente possui \textbf{uma} notável contribuição para o desenvolvimento das competências específicas da área de Ciências Humanas e Sociais Aplicadas, torna -se necessário direcionar o olhar para essas competências e analisar o percurso educacional.
No exame geral dessas competências, evidenciam -se a análise, a \textbf{contextualização}, o reconhecimento e a participação dos estudantes, sujeitos protagonistas do processo de aprendizagem.
Cabe ao professor a mobilização de estratégias diversificadas para o acompanhamento das aprendizagens, por meio de diversificadas maneiras : textual, oral, imagética, discursiva, entre outras.
Dado o caráter prático da Sociologia, o componente prepara o estudante para a \textbf{problematização} da realidade e ação sobre a mesma, algo \textbf{que} está evidenciado na proposta da área de Ciências Humanas e Sociais Aplicadas na BNCC.
Nesse sentido, a avaliação da aprendizagem vai muito além da \textbf{mera} verificação de conteúdos, direcionando -se à reflexão sobre o percurso formativo, considerando o desenvolvimento de competências e habilidades, ou seja, o conjunto de conhecimentos e as ações, somadas às tomadas de decisão feitas pelos estudantes, \textbf{ancoradas} pelo saber escolar.

% matched lemmas: abordagem, ademais, amparar, ancorar, aprimorar, aproximar, aqui, argumentar, assimilar, atentar, ator, atual, autor, basear, cartográfico, categoria, compromisso, conceitual, condizente, constitutivo, contemporâneo, contextualização, contudo, depender, dignidade, direcionamento, dizer, embasar, encaminhamento, enquanto, ensino-aprendizagem, equidade, espírito, exigir, exposição, exposto, face, fronteira, fundamento, hegemônico, historicamente, homem, horizonte, imediato, imprescindível, inerente, integralidade, inter-relação, interessante, memorização, mero, metodológico, método, norteador, novidade, operar, percepção, permear, preciso, problematização, prova, que, raciocínio, restringir, salientar, sistematizar, sujeito, tarefa, teoria, termos, teórico, um
\end{textsample}
